% !TeX root = ../main.tex
% ============================================================
%  sections/07_discussion.tex
% ============================================================

\section{Discussion}
\label{sec:discussion}

The experiments in \Cref{sec:experiments} support three substantive
conclusions about taxing closed-loop supply chains: the emissions-minimising
tax is instance-specific and sits on a saturating frontier; cleaner production
and remanufacturing technology, not transportation, is the binding abatement
lever; and emissions are governed by demand and by production and upstream
emission intensities. We discuss what each implies for managers and
regulators, then state the limits of the analysis.

\subsection{The cost--emission trade-off and the saturating frontier}
\label{sec:disc:tradeoff}

The optimal taxes recovered by the bilevel model vary by an order of magnitude
across network configurations ($\tau^{\star}\in[0.02,0.20]$), and the
follower's binary facility, technology, capacity, and transport decisions make
its response change in discrete blocks rather than smoothly. The most
consequential feature is the \emph{saturation} of the emission--tax response:
beyond a moderate rate, further taxation buys almost no additional abatement,
so the emissions-minimising tax is attained over a wide, near-flat band
(\Cref{sec:exp:tradeoff}). A regulator that optimises emissions alone is
therefore ill-posed in this regime --- the objective alone cannot distinguish a
moderate tax from a punitive one that achieves a rounding-error reduction. The
lexicographic least-rate criterion is what restores well-posedness, and the
practical lesson is that an emissions target must be paired with a cost
criterion or a minimum meaningful-abatement threshold; otherwise the
``optimal'' tax is an artifact of where the grid happens to end.

\subsection{Robustness of the qualitative findings}
\label{sec:disc:robustness}

Two features make the findings more than instance anecdotes. The factor
rankings are \emph{stable across network sizes} --- demand is the leading cost
driver in three of four screened instances, and technology dominates abatement
in all four scenario instances --- so the qualitative conclusions are not
artifacts of a single configuration. And the underlying solves are
trustworthy: every one of the 95 committed runs is validated, the largest
relative gap is $1.9\times10^{-5}$, and the procedure converges in two to
three iterations. The robust claim is thus not that one tax rate is universal,
but that \emph{which lever abates and which inputs matter} is consistent across
the tested networks.

\subsection{Managerial implications}
\label{sec:disc:managerial}

For supply-chain managers, the clearest signal is where to invest abatement
effort. Expanding the production and remanufacturing technology menu reduces
optimal emissions by $26$--$33\%$, while expanding the transport-mode
catalogue is effectively inert; on the larger I25 networks the richer
technology set is in fact jointly cheaper and cleaner. Capital and managerial
attention are therefore better directed at cleaner production and
remanufacturing capability than at modal shift in these networks. The
sensitivity analysis sharpens the point: demand and the production and
upstream (scope-3) emission intensities --- \emph{what is made and procured} ---
dominate both cost and emissions, so demand exposure and supplier emission
intensity are the key risks to manage, whereas trucking and storage parameters
are second-order.

\subsection{Policy implications}
\label{sec:disc:policy}

For regulators, the instance-specific variation in $\tau^{\star}$ argues
against a one-size-fits-all carbon price: if sectoral network structure is
known, the bilevel model supports differentiated calibration, and if it is
not, it still provides a reproducible way to test representative networks
before setting a rate. The saturating frontier carries a second, cautionary
message: once a tax is high enough to trigger the available abatement, raising
it further imposes cost on the supply chain for negligible environmental
return. A tax instrument is most effective when calibrated to the supply
chain's abatement
\emph{response}, which is exactly what the bilevel formulation makes
observable; pushing past that point is economically wasteful. Because the
binding abatement lever here is production/remanufacturing technology, a carbon
tax may also be more effective when paired with instruments that lower the cost
of adopting cleaner technology.

\subsection{Limitations of the present study}
\label{sec:disc:limitations}

Several modelling and computational choices bound the analysis.

\paragraph{Generated deterministic data.}
Demand, return fractions, costs, emissions, and capacities are generated and
deterministic. The results are reproducible but are not calibrated to a
specific industrial case, so the magnitudes (\eg the $26$--$33\%$ technology
effect) should be read as model-internal rather than as field estimates.

\paragraph{Lexicographic and optimistic tie-breaking.}
The leader minimises emissions and then the tax, and the follower breaks cost
ties toward lower emissions. These optimistic conventions favour
lower-emission outcomes; a pessimistic regulator that cannot assume favourable
follower behaviour would face a less forgiving problem, which we do not model here.

\paragraph{Bounded benchmark and horizon.}
The study covers the eight T3--T6 base instances; the four T12 instances are
intractable within the compute budget, so the longest horizons are out of
reach without further algorithmic work. The sensitivity screen covers four
bases and the Morris design covers one, so the global screening is preliminary.

\paragraph{Structural observations.}
Warehouses never open in any optimal response under the generated data, so the
warehouse echelon is effectively dormant in these instances --- a property of
the data and cost structure rather than evidence that storage is never useful.
We also do not decompose abatement into per-facility or per-technology
contributions; the scenario design isolates the technology \emph{lever} but not
the individual mechanisms within it.

\paragraph{Single instrument and no external benchmark.}
The leader wields one discrete tax index over $\{0,\ldots,\taxub\}$; caps,
subsidies, cap-and-trade, and dynamic trajectories are outside the present
scope. We also do not benchmark the \DBP against an alternative exact method
(\eg a KKT reformulation or exhaustive grid beyond the enumeration
cross-check), so we make no cross-method performance claim.
